\documentclass[../../main.tex]{subfiles}

\begin{document}

\section{无限集与选择公理}\label{section:ch01-无限集与选择公理}

我们已经得到了判定集合为无限集的一些法则。例如，若集合 $A$ 有一个可数无限子集，或者存在 $A$ 与其真子集之间的一一对应，则 $A$ 必定为无限集。我们将要证明，这些性质中的任何一个都足以刻画无限集。而这个证明又把我们引向讨论目前我们还没有提到过的逻辑学中的一个问题——选择公理。

\begin{theorem}[无限集的刻画]\label{theorem:ch01-无限集刻画}
设 $A$ 是一个集合。关于 $A$ 的下列条件等价：
\begin{enumerate}[(1)]
\item 存在一个单射 $f: \mathbb{Z}_{+} \longrightarrow A$.
\item $A$ 与其真子集之间存在一个一一对应.
\item $A$ 是一个无限集.
\end{enumerate}
\end{theorem}

\par
\begin{figure}[htbp]
\centering
\includegraphics[width=0.6\textwidth]{ch01-fig9-1}
\caption{将可数集 $B$ 右移一位，$A - B$ 上恒等}\label{figure:ch01-无限集映射}
\end{figure}

\begin{proof}
我们证明蕴涵关系 (1) $\Longrightarrow$ (2) $\Longrightarrow$ (3) $\Longrightarrow$ (1)。为了证明 (1) $\Longrightarrow$ (2)，设存在一个单射 $f: \mathbb{Z}_{+} \longrightarrow A$。将像集 $f(\mathbb{Z}_{+})$ 记为 $B$，将 $f(n)$ 记为 $a_{n}$。由于 $f$ 是一个单射，当 $n \neq m$ 时，$a_{n} \neq a_{m}$。用公式
\begin{align*}
g(a_{n}) = a_{n+1}, \quad \text{对于 } a_{n} \in B,
\end{align*}
\begin{align*}
g(x) = x, \quad \text{对于 } x \in A - B
\end{align*}
定义一个映射
\begin{align*}
g: A \longrightarrow A - \{a_{1}\}.
\end{align*}
这个映射 $g$ 如\cref{figure:ch01-无限集映射}所示，容易验证它是一个一一对应。


蕴涵关系 (2) $\Longrightarrow$ (3) 正好是\cref{corollary:ch01-有限不自映射}的逆否命题，我们已经证明过了。为了证明 (3) $\Longrightarrow$ (1)，假定 $A$ 为无限集，并且“用归纳法”定义一个单射 $f: \mathbb{Z}_{+} \longrightarrow A$。

首先，因为集合 $A$ 非空，可在 $A$ 中取一点 $a_{1}$，定义 $f(1)$ 为所选取的那个点。

其次，假定 $f(1), \cdots, f(n-1)$ 已经有定义，我们来定义 $f(n)$。集合 $A - f(\{1, \cdots, n-1\})$ 是非空的，因为如果它是空集的话，映射 $f: \{1, \cdots, n-1\} \longrightarrow A$ 就是满射，由此推出 $A$ 是有限的。因此，我们可以在集合 $A - f(\{1, \cdots, n-1\})$ 中取一个元素，并且就定义 $f(n)$ 为这个元素。应用“归纳原理”，对于所有 $n \in \mathbb{Z}_{+}$，$f$ 有定义。

容易看出 $f$ 是单射。如果假定 $m < n$，则 $f(m)$ 属于集合 $f(\{1, \cdots, n-1\})$，而根据定义 $f(n)$ 不属于这个集合。所以 $f(n) \neq f(m)$。
\end{proof}

让我们更加仔细地改进这种“归纳”证明，以便我们对归纳定义原理的使用更加清楚。

给定一个无限集 $A$，我们打算用下式
\begin{align}
f(1) = a_{1},\label{align:ch01-选择归纳}
\end{align}
\begin{align*}
f(i) = [A - f(\{1, \cdots, i-1\})] \text{ 的一个任意元素，对于 } i > 1,
\end{align*}
归纳地定义 $f: \mathbb{Z}_{+} \longrightarrow A$。但这并不是一个理想的归纳公式！因为 $f(i)$ 不是由 $f \mid \{1, \cdots, i-1\}$ 唯一确定的。

在这方面，这个公式与我们在证明\cref{lemma:ch01-正整数无限子集}时考虑过的归纳公式完全不同。那里有一个 $\mathbb{Z}_{+}$ 的无限子集 $C$，并且用下式
\begin{align*}
h(1) = C \text{ 的最小元},
\end{align*}
\begin{align*}
h(i) = [C - h(\{1, \cdots, i-1\})] \text{ 的最小元}, \quad \text{对于 } i > 1,
\end{align*}
定义 $h$，而这个公式中 $h(i)$ 由 $h \mid \{1, \cdots, i-1\}$ 唯一确定。

我们说\eqref{align:ch01-选择归纳}不是一个理想的归纳公式的另一个原因是，因为如果不是这样，那么由归纳定义原理得出满足\eqref{align:ch01-选择归纳}的唯一的一个函数 $f: \mathbb{Z}_{+} \longrightarrow A$。但是很难想像可以由\eqref{align:ch01-选择归纳}唯一确定 $f$。事实上，$f$ 的这个“定义”含有无限多种选择。

我们必须指出，\cref{theorem:ch01-无限集刻画}中给出的证明，实际上不是一个证明。事实上，我们在迄今所讨论过的集合论的性质的基础上，还不可能证明这个定理。必须增加一些东西才行。

我们在前面讲过描述指定集合的一些可行的方法：
\begin{enumerate}[(1)]
\item 列出它的元素以定义集合，或者取一个给定的集合 $A$，然后按照元素是否满足某给定的性质而确定它的一个子集 $B$.
\item 取给定的集族元素的并与交，或者两个集合的差.
\item 取给定集合的所有子集的集合.
\item 取集合的笛卡儿积.
\end{enumerate}

现在函数 $f$ 的法则实际上是一个集合，即 $\mathbb{Z}_{+} \times A$ 的子集。因此，要证明函数 $f$ 的存在性，就必须用可允许的组成集合的方法，做出 $\mathbb{Z}_{+} \times A$ 的适当的子集。但是已经给定的这些方法却达不到这一目的。我们需要有断定集合存在性的新方法。为此要在所有可允许的组成集合的方法中加进下述公理：

\begin{axiom}[选择公理]\label{axiom:ch01-选择公理}
给定由两两无交的非空集合构成的一个族 $\mathcal{A}$，存在一个集合 $C$，使得 $C$ 与 $\mathcal{A}$ 的每一个元素恰好有一个公共元，即对于每一个 $A \in \mathcal{A}$，集合 $C \cap A$ 包含着唯一的一个元素。
\end{axiom}

选择公理所述存在的那个集合 $C$，可以将其设想成从 $\mathcal{A}$ 的每一个集合 $A$ 中选取一个元素而得到。

选择公理看起来显然是一个十分清楚的论断。事实上，大多数数学家现在都把它作为数学基础中集合论的一部分。但是在过去的若干年里，围绕这个与集合论有关的特殊论断有过激烈的争论。一些数学家对于通过它所证明的一些定理只是勉强承认。其中的一个就是我们马上要讨论的良序定理。现在就可以直接使用选择公理来克服前面证明中的困难。我们首先证明选择公理的一个简单推论。

\begin{lemma}[选择函数的存在性]\label{lemma:ch01-选择函数}
给定非空集合的一个族 $\mathcal{B}$（未必是两两无交的）。则存在一个函数
\begin{align*}
c: \mathcal{B} \longrightarrow \bigcup_{B \in \mathcal{B}} B,
\end{align*}
使得对于每一个 $B \in \mathcal{B}$，$c(B)$ 是 $B$ 的一个元素。

函数 $c$ 称为族 $\mathcal{B}$ 的一个选择函数（choice function）。
\end{lemma}

\begin{proof}
这个引理与选择公理之间的区别是：引理中族 $\mathcal{B}$ 中的集合不要求是两两无交的。例如，可以允许 $\mathcal{B}$ 为一个给定集合的非空子集族。

对于 $\mathcal{B}$ 的一个元素 $B$，定义集合 $B'$ 为：
\begin{align*}
B' = \{(B, x) \mid x \in B\}.
\end{align*}
即 $B'$ 是所有有序偶对的族，有序偶对的第一个坐标是集合 $B$，第二坐标是 $B$ 的一个元素。$B'$ 是笛卡儿积
\begin{align*}
\mathcal{B} \times \bigcup_{B \in \mathcal{B}} B
\end{align*}
的一个子集。因为 $B$ 至少含有一个元素 $x$，那么 $B'$ 至少包含元素 $(B, x)$，因此 $B'$ 非空。

现在证明：如果 $B_{1}$ 和 $B_{2}$ 是 $\mathcal{B}$ 中两个不同的集合，则 $B_{1}'$ 和 $B_{2}'$ 无交。因为 $B_{1}'$ 的元素形如 $(B_{1}, x_{1})$，$B_{2}'$ 的元素形如 $(B_{2}, x_{2})$。两者的第一个坐标不同，所以这样的两个元素不会相等。

我们构造一个族
\begin{align*}
\mathcal{C} = \{B' \mid B \in \mathcal{B}\};
\end{align*}
它是
\begin{align*}
\mathcal{B} \times \bigcup_{B \in \mathcal{B}} B
\end{align*}
的无交非空子集的一个族，由\cref{axiom:ch01-选择公理}知，存在一个集合 $c$，它与 $\mathcal{C}$ 的每一个元素恰好有一个公共元。我们要证明 $c$ 就是所需要的选择函数的法则。

首先，$c$ 是
\begin{align*}
\mathcal{B} \times \bigcup_{B \in \mathcal{B}} B
\end{align*}
的一个子集。其次，$c$ 恰好含有每一个集合 $B'$ 的一个元素。因此，对于每一个 $B \in \mathcal{B}$，集合 $c$ 恰好含有第一个坐标为 $B$ 的一个有序偶对 $(B, x)$。所以 $c$ 事实上就是从族 $\mathcal{B}$ 到集合 $\bigcup_{B \in \mathcal{B}} B$ 的一个函数的法则。最后，若 $(B, x) \in c$，则 $x$ 属于 $B$，于是 $c(B) \in B$，证毕。
\end{proof}

\cref{theorem:ch01-无限集刻画}的第二个证明 应用上面这个引理，可以更加严格地证明\cref{theorem:ch01-无限集刻画}。给定一个无限集 $A$，我们希望做出一个单射 $f: \mathbb{Z}_{+} \longrightarrow A$。记 $A$ 的所有非空子集的族为 $\mathcal{B}$。刚才证明的\cref{lemma:ch01-选择函数}保证了 $\mathcal{B}$ 的选择函数的存在性，即有一个函数
\begin{align*}
c: \mathcal{B} \longrightarrow \bigcup_{B \in \mathcal{B}} B = A,
\end{align*}
使得对于每一个 $B \in \mathcal{B}$，$c(B) \in B$。用归纳公式
\begin{align}
f(1) = c(A),\label{align:ch01-严格归纳}
\end{align}
\begin{align*}
f(i) = c(A - f(\{1, \cdots, i-1\})), \text{对于 } i > 1,
\end{align*}
定义一个函数 $f: \mathbb{Z}_{+} \longrightarrow A$。因为 $A$ 是无限集，集合 $A - f(\{1, \cdots, i-1\})$ 非空。因此，上式右边有意义。因为这个公式是用 $f \mid \{1, \cdots, i-1\}$ 唯一定义的 $f(i)$，所以能用\cref{theorem:ch01-归纳定义原理}。我们得到了对于所有 $i \in \mathbb{Z}_{+}$ 满足\eqref{align:ch01-严格归纳}的唯一函数 $f: \mathbb{Z}_{+} \longrightarrow A$。$f$ 的单射性的证明如前。

为了给出\cref{theorem:ch01-无限集刻画}合乎逻辑的证明，要用到选择函数。当我们强调了这一点之后，还必须回过头来说清楚大多数数学家并不是这样做的。他们毫不迟疑地给出如同第一种情形那样的证明，其中包含有无穷多次的任意选择。他们知道自己确实用了选择公理。他们也知道如果需要的话，那就可以通过引进一个特殊的选择函数，而使其证明达到逻辑上更令人满意的程度。但是他们常常不这样做。

我们也是如此。你将在本书中很少找到选择函数的特定用处。只是在证明变得含混时才引进一个选择函数，但是在许多证明中，做了无穷多次任意选择，每当这种情形，我们实际上暗中使用了选择公理。

我们必须承认，在本书上一节的一个证明中，从无穷多种可供选择的函数中作出了一个确定的函数，似乎没有用到选择公理。这种误解要加以澄清。请你找找看，我们说的究竟是哪一个证明。

为选择公理加一个最后的附注。这个公理有两种形式。一种称为有限选择公理（finite axiom of choice），它说对于一个两两无交的非空集合的有限族 $\mathcal{A}$，存在一个集合 $C$，恰好与 $\mathcal{A}$ 的每一个元素只有一个公共元。我们一直需要选择公理的这种弱形式，以前我们一直自由地使用它而不加说明。对于有限选择公理，任何一个数学家也不会觉得不妥，它是所有集合论的组成部分。换言之，一个证明如果只涉及有限多个任意选择时，没有人会提出疑问。

选择公理的强形式是对非空集合任意的一个族 $\mathcal{A}$ 而言的，只有它才叫做“选择公理”。每当数学家写道：“这个证明依赖于选择公理”时，它总是指选择公理的这种强形式。

\subsection*{习题}

\begin{exercise}\label{exercise:ch01-选择公理-1}
不用选择公理定义一个单射 $f: \mathbb{Z}_{+} \longrightarrow X^{\omega}$，其中 $X$ 为二元素集 $\{0, 1\}$。
\end{exercise}

\begin{exercise}\label{exercise:ch01-选择公理-2}
如果有可能的话，试对下列各个集族不用选择公理而求出来一个选择函数。
\begin{enumerate}[(a)]
\item $\mathbb{Z}_{+}$ 的所有非空子集的族 $\mathcal{A}$.
\item $\mathbb{Z}$ 的所有非空子集的族 $\mathcal{B}$.
\item 有理数集 $\mathbb{Q}$ 的所有非空子集的族 $\mathcal{C}$.
\item $X^{\omega}$ 的所有非空子集的族 $\mathcal{D}$，其中 $X = \{0, 1\}$.
\end{enumerate}
\end{exercise}

\begin{exercise}\label{exercise:ch01-选择公理-3}
设 $A$ 是一个集合，$\{f_{n}\}_{n \in \mathbb{Z}_{+}}$ 是加标单射的一个族，其中
\begin{align*}
f_{n}: \{1, \cdots, n\} \longrightarrow A.
\end{align*}
证明 $A$ 是一个无限集。你能不用选择公理定义一个单射 $f: \mathbb{Z}_{+} \longrightarrow A$ 吗？
\end{exercise}

\begin{exercise}\label{exercise:ch01-选择公理-4}
在第\cref{section:ch01-可数集与不可数集}节中，有一个定理的证明包含有无穷多次任意选择，它是哪一个呢？重新写出一个证明，以便弄清楚选择公理的用法。（前面有一些习题也用到过选择公理。）
\end{exercise}

\begin{exercise}\label{exercise:ch01-选择公理-5}
(a) 应用选择公理证明：若 $f: A \longrightarrow B$ 为满射，则 $f$ 有一个右逆 $h: B \longrightarrow A$。

(b) 证明：若 $f: A \longrightarrow B$ 为单射，$A$ 非空，则 $f$ 有一个左逆。这需要用选择公理吗？
\end{exercise}

\begin{exercise}\label{exercise:ch01-选择公理-6}
朴素集合论中大多数著名的悖论，都以这种或那种方式与“所有集合的集合”这个概念有关。前面曾经给出过的用来构成集合的各种法则，没有一个允许我们考虑这样的集合。我们有足够的理由指出，这个概念本身是自相矛盾的。假定用 $\mathcal{A}$ 表示“所有集合的集合”。
\begin{enumerate}[(a)]
\item 证明 $\mathcal{P}(\mathcal{A}) \subset \mathcal{A}$，并且由此引出矛盾.
\item （Russell 悖论）设 $\mathcal{B}$ 是 $\mathcal{A}$ 的一个子集，它由所有不是它自身的元素的集合组成，即
\begin{align*}
\mathcal{B} = \{A \mid A \in \mathcal{A} \text{ 并且 } A \notin A\}.
\end{align*}
（当然，可能没有集合 $A$ 会使得 $A \in A$。这时，$\mathcal{B} = \mathcal{A}$。）$\mathcal{B}$ 是不是它自身的一个元素呢？
\end{enumerate}
\end{exercise}

\begin{exercise}\label{exercise:ch01-选择公理-7}
设 $A$ 和 $B$ 是两个非空集合。如果从 $B$ 到 $A$ 中有一个单射，而从 $A$ 到 $B$ 中没有单射，我们说 $A$ 有比 $B$ 较大的基数（greater cardinality）。
\begin{enumerate}[(a)]
\item 根据\cref{theorem:ch01-无限集刻画}得出，每一个不可数集基数都大于 $\mathbb{Z}_{+}$ 的基数.
\item 证明：若 $A$ 的基数大于 $B$ 的基数，$B$ 的基数大于 $C$ 的基数，则 $A$ 的基数大于 $C$ 的基数.
\item 求出无限集的一个序列 $A_{1}, A_{2}, \cdots$，使得对于每一个 $n \in \mathbb{Z}_{+}$，集合 $A_{n+1}$ 的基数大于 $A_{n}$ 的基数.
\item 求出一个集合，使得对于任意 $n$，它的基数大于 $A_{n}$ 的基数.
\end{enumerate}
\end{exercise}

\begin{exercise}\label{exercise:ch01-选择公理-8}
证明 $\mathcal{P}(\mathbb{Z}_{+})$ 与 $\mathbb{R}$ 的基数相同。[提示：你可以利用这样一个事实，即每一个实数有唯一的一个十进制展开，只要在这个展开式中不允许出现无穷多个 $9$ 构成的串。]

集合论中一个著名的猜想，称为连续统假设，它断言不存在基数大于 $\mathbb{Z}_{+}$ 小于 $\mathbb{R}$ 的集合。广义连续统假设则断言，对于无限集 $A$，不存在基数大于 $A$ 小于 $\mathcal{P}(A)$ 的集合。令人吃惊的是，这两个论断与集合论的常用公理是独立的。比较易读的阐述可参见 $[\mathrm{Sm}]$。
\end{exercise}

\end{document}